% =====================================================================
%  RESUMEN Y PALABRAS CLAVE
%  Responsable: TODOS (se redacta AL FINAL, cuando el informe esté cerrado)
%  Extensión objetivo: 250 a 300 palabras
%
%  En este formato el resumen NO es un capítulo numerado: va en el entorno
%  abstract de IEEEtran, justo debajo del bloque de autores y a una sola
%  columna. No lleva \chapter ni \section.
% =====================================================================

\begin{abstract}
% TODO: Un solo párrafo compacto que responda, en este orden:
%   1) Qué caso se analiza  -> Mina Marlin, Montana Exploradora / Goldcorp,
%      municipios de San Miguel Ixtahuacán y Sipacapa, San Marcos, Guatemala;
%      oro y plata; operó de 2005 a 2017.
%   2) Qué auditoría se estudia -> la Human Rights Assessment of Goldcorp's
%      Marlin Mine, de On Common Ground Consultants (2010), evaluación
%      socioambiental independiente de más de 200 páginas.
%   3) Qué metodología se aplicó en ESTE informe -> revisión documental de
%      fuentes primarias, análisis bajo ISO 19011:2018 y clasificación de
%      hallazgos en no conformidades.
%   4) Hallazgos centrales -> los siete ejes; en especial la ausencia de
%      consulta previa (Convenio 169 de la OIT), los riesgos de drenaje
%      ácido de roca y calidad de agua, la seguridad privada armada y la
%      debilidad de los mecanismos de queja.
%   5) El resultado propuesto -> plan de mejora continua bajo el ciclo PHVA
%      alineado a los IFC Performance Standards, el GISTM e IRMA.
%
% REGLA: el resumen NO lleva citas, ni tablas, ni figuras. Solo prosa.

\pendiente{redactar el resumen (250--300 palabras)}
\end{abstract}

\begin{IEEEkeywords}
% TODO: de 4 a 6 palabras clave separadas por punto y coma.
% Sugeridas: auditoría socioambiental; Mina Marlin; ISO 19011; ciclo PHVA;
% derechos humanos; no conformidad.
\pendiente{4 a 6 palabras clave}
\end{IEEEkeywords}
